Sea $\Psi\in H_2$. Por el punto anterior, las cuatro funciones $e^{i(nt+ms)}$ forman una base de $H_2$, y el isomorfismo $T$ identifica cada una con un elemento $\ket a\otimes\ket b$ de la base de $H_1\otimes H_1$, con $n=a+1$ y $m=b+1$. Por eso existen números complejos únicos $c_{ab}$ tales que
\[
  \Psi=\sum_{a,b\in\{0,1\}}c_{ab}\,\ket a\otimes\ket b ,
\]
que, en términos de funciones, significa
\[
  \Psi(t,s)=c_{00}\,e^{i(t+s)}+c_{01}\,e^{i(t+2s)}+c_{10}\,e^{i(2t+s)}+c_{11}\,e^{2i(t+s)} .
\]
Los coeficientes se ordenan en la matriz
\[
  C=\begin{pmatrix}c_{00}&c_{01}\\ c_{10}&c_{11}\end{pmatrix}.
\]

\paragraph{Caracterización.} La onda $\Psi$ es de la forma $\psi(t)\varphi(s)$ con $\psi,\varphi\in H_1$ si y solo si
\[
  \det C=c_{00}c_{11}-c_{01}c_{10}=0 .
\]
En la notación del enunciado, con $\Psi=\sum c_{nm}e^{i(nt+ms)}$, la condición es $c_{11}c_{22}=c_{12}c_{21}$.

\paragraph{Demostración.} Si $\psi=x_0\ket0+x_1\ket1$ y $\varphi=y_0\ket0+y_1\ket1$, entonces
\[
  \psi(t)\varphi(s)=\sum_{a,b}x_a y_b\,\ket a(t)\,\ket b(s),
\]
así que $c_{ab}=x_a y_b$. En particular $c_{00}c_{11}=x_0y_0x_1y_1=c_{01}c_{10}$, es decir, $\det C=0$.

Recíprocamente, supóngase que $\det C=0$. Si $C=0$, entonces $\Psi=0$, que es el producto $0\cdot0$. Si $C\neq0$, alguna fila de $C$ es no nula; sea $y=(y_0,y_1)$ esa fila. Como el determinante es cero, la otra fila es un múltiplo de $y$: existen números $x_0,x_1$ tales que $(c_{a0},c_{a1})=x_a\,y$ para $a=0,1$. Esto significa $c_{ab}=x_a y_b$, y por tanto
\[
  \Psi(t,s)=\psi(t)\varphi(s), \qquad \psi=\sum_a x_a\ket a,\quad \varphi=\sum_b y_b\ket b .
\]

\paragraph{Relación con la probabilidad.} La condición es la misma que aparece con dos variables aleatorias. Si $p_{ab}$ es la distribución conjunta de dos variables que toman los valores $0$ y $1$, las variables son independientes cuando la conjunta es el producto de sus marginales, $p_{ab}=p_a\,q_b$, y hay distribuciones conjuntas que no son el producto de sus marginales. Para una tabla de $2\times2$, además, la independencia equivale a $p_{00}p_{11}=p_{01}p_{10}$, que es la condición $\det C=0$ escrita con probabilidades. Del mismo modo, una onda de $H_2$ es un producto $\psi(t)\varphi(s)$ cuando $c_{ab}=x_ay_b$, y las ondas de $H_2$ que no son de esa forma son los estados entrelazados.

\paragraph{Ejemplo.} La onda $\Psi=\frac1{\sqrt2}\bigl(\ket0\otimes\ket0+\ket1\otimes\ket1\bigr)$, que aparece en el punto siguiente, tiene $C=\frac1{\sqrt2}\begin{pmatrix}1&0\\0&1\end{pmatrix}$ y $\det C=\frac12\neq0$. Por tanto no es de la forma $\psi(t)\varphi(s)$: es un estado entrelazado.
